\section{Conclusion}
\label{sec:conclusion}

This paper reported the design, firmware implementation, an honest
object-oriented architecture audit, and a hardware-validated correctness
campaign for Steistat, our own open-source, low-cost AD5941/RP2040
potentiostat. Direct, file-by-file and register-by-register comparison
against Analog Devices' own reference examples, together with the
instrument's own diagnostic instrumentation, found and fixed eleven
concrete firmware defects: an excitation/sensing-loop misconfiguration, an
ADC zero-point convention error, a calibration-resistor declaration that
scaled every reported current by a fixed, linearity-invisible factor, a
runtime calibration command that silently failed to reach the measurement
path, a transimpedance-gain calibration frozen at power-up, an ADC
programmable-gain enum mismatch, an inverted bias-potential sign
convention, a switch-matrix misconfiguration that partly sank cell current
into the wrong node, an AFE left hibernating mid-measurement, and two
unbounded wait loops capable of wedging the instrument indefinitely.
Several of these existed in three or four verbatim copies, once per
affected method class -- the direct, measured cost of the duplication our
own source-level audit had flagged -- and each had to be found and fixed
that many times over.

On the corrected firmware, this paper's own fresh hardware measurements
show cyclic voltammetry is quantitatively accurate: against a traceable
Randles dummy cell, five replicate sweeps recover
$R=\SI{10605.1\pm0.9}{\ohm}$ against a \SI{10560}{\ohm} nominal value
($+0.427\,\%$ error, $\%\mathrm{RSD}=0.0084\,\%$, $R^2\geq0.9999986$ on
every sweep), independently reproducing an earlier $n=10$ measurement on
this same corrected signal chain to within \SI{0.1}{\ohm}; against a plain
\SI{10}{\kilo\ohm} resistor, the same protocol recovers \SI{10005.9}{\ohm}
($+0.06\,\%$ error, $R^2=0.9999988$). Chronoamperometry's initial,
step-onset current is correctly signed and monotonic across five applied
potential steps, and square-wave and differential-pulse voltammetry
produce smooth, bounded differential-current curves free of the
railing/noise previously reported. This paper also reports, honestly and
without correcting it here, a real defect found in the same diagnostic
work: the \emph{sustained} current reported by CA, SWV, and DPV does not
track the applied potential, evidence points to a cause upstream of the
sensing and conversion stage rather than within it, and root-causing this
is left to future work rather than folded into an overclaimed validation.
An accompanying, evidence-based audit of the firmware's object-oriented
design found real but limited object orientation, and this paper closes
every gap that audit raised: all six method classes, including cyclic
voltammetry (previously outside the class structure entirely, now wrapped
in a new \texttt{C\_CV} class), sit under one polymorphic
\texttt{C\_ElectrochemicalMethod} base with a genuine inheritance
relationship, and the central parameter store, \texttt{C\_DataStorage}, is
fully field-encapsulated rather than an all-public-fields struct -- both
changes verified on physical hardware to leave this paper's own CV
accuracy result unchanged. This is not a claim of a maximally-abstracted
architecture: one inheritance level, no call-site dispatch through
base-class pointers, and mechanical get/set accessors are what this paper
delivers, not a richer domain design. Future work should characterize the
fitted RCAL1 resistor against a traceable standard, root-cause and fix the
sustained CA/SWV/DPV defect, re-validate EIS, and complete the
sample-timing and on-chip-sequencer improvements identified in
Section~\ref{sec:limitations}.
